\section{问题重述}

\subsection{问题背景}

大语言模型通过预训练从大量文本中学习通用能力。经典标度律以模型参数量 $N$ 和训练 token 数 $D$ 为自变量，用幂律形式刻画验证集交叉熵损失 $L_0(N,D)=E+AN^{-\alpha}+BD^{-\beta}$ 的下降\cite{hoffmann,bahri}。扩大 $N$ 或 $D$ 都需要算力；按本题采用的近似，基础训练开销 $C\approx6ND$，因而有限预算下必须权衡模型规模与数据量。

仅比较 $N,D,C$ 仍不足以描述真实训练：不同语料的质量和领域组成会改变同等投入的效果。因而需要先建立数据质量 $Q$ 的可比口径与领域配比 $\bp$ 的损失关系，再研究如何把二者接入经典标度律，形成广义标度律 $L(N,D,Q,\bp)$。这一递进把数据选择、训练成本以及最终能力评测纳入同一资源配置问题。赛题的附件 A 提供质量信号与配比实验，附件 B 提供标度律相关数据，附件 C 提供模型、算力与能力评测信息。

\subsection{问题描述}

本题依次要求完成四项数学任务：评价训练语料的质量并刻画领域配比的损失效应；在经典标度律中引入质量与配比，建立可检验的广义标度律；在非线性算力成本约束下优化训练配置；把规模和非规模因素映射到能力前沿，作有条件的演进预测。前一问的输出是后一问的输入，数据口径和检验边界也随之累积。

\textbf{问题一：数据质量与领域配比。}建模对象是文档、领域和混合语料的质量，以及 17 个训练领域的配比对 13 个验证域损失的影响；约束是质量信号方向和量纲不一、部分领域缺少直接质量映射，且配比份额非负、总和为 1。用 A1 拟合文档质量评分与预处理规则，冻结后在扩展集 A2/A3 上检查分数与冲突诊断的稳定性；用 A16 的领域映射和 A18 的文本特征把质量扩展到 17 域。A4/A5 按索引连接后拟合配比—损失关系；A6/A7 作同尺度留出检验，A8–A11 作跨尺度直接检验；A12–A15 仅用于估算或外推情景讨论，不当作真实训练验证。

\textbf{问题二：广义标度律。}建模对象是损失对 $N,D,Q,\bp$ 的依赖；约束是理想质量与基准配比下应回到经典基线，A、B 两组损失的评测管线和质量尺度不能直接等同，且配比始终满足份额约束。以 B1 拟合经典基线，以 B6 拟合质量进入损失的候选形式，并仅用 B7 相对 B6 新增的配置作独立留出检验；B2/B3 作半合成或插值参照，B4/B5 作不同来源的趋势对照，B8 作方向冲突压力检查；B9/B10 按模型名及 $N,D$ 对齐，B10 仅作估算对照，B11/B12 拟用于来源审计。问题一的领域质量和配比关系经固定接口引入；其中配比通道在满足训练与留出条件前只作为探索性扩展。

\textbf{问题三：算力约束下的训练优化。}建模对象是在给定预算 $C$ 下的参数量 $N$、数据量 $D$ 与质量 $Q$ 的最优配置；约束包括训练、提质和长上下文三项成本，质量下界由问题一的推荐配比确定，主方案固定该配比。读取问题一的质量与配比接口，并在问题二完成留出验收后使用其冻结的标度律参数，不在本问重新拟合；三类提质成本函数来自赛题，C7 提供可行的上下文长度用于敏感性分析。比较质量下界、内点和上界解，判断预算变化时是否存在以及何处出现结构性转移。

\textbf{问题四：能力演进与预测。}建模对象是开源模型能力前沿及其规模与非规模变化；约束是能力、时间、模型类型和开放权重口径必须一致，且损失与 Benchmark 分数不在同一尺度。C1 提供榜单能力，C2 在同批记录上补充发布日期；C4 补充规模和算力并与 C2 匹配；C8 用于逐任务聚合和饱和度检查；C6 分层拟合损失到得分的桥接，C3 只作历史混合口径的敏感性对照。结合前问的标度律与预算情景建立可解释的演进模型，对历史前沿作分解与回测，再讨论算力增速放缓时 12/24 个月的条件预测及其不确定性。

\section{问题分析}

\subsection{总体思路}

四问按“测量数据—描述损失—优化训练—解释能力”递进（图~\ref{fig:flow}）。问题一从文档质量 $Q(x)$ 聚合到领域质量 $Q_d$，再按配比得到语料质量 $Q(\bp)$；问题二在经典的 $L_0(N,D)$ 上引入 $Q$，并探索配比作用。经典律只是广义律在理想质量、基准配比下的一个切片。问题三以标度律为目标，新增预算 $C$、提质成本与上下文长度约束，求得最优配置；问题四再引入时间 $t$ 和 Benchmark 度量，把损失变化转为能力前沿的解释与预测。各问的冻结结果通过 P1–P4 接口传递，未通过留出检验的配比通道不进入下游主优化。

数据使用保持\textbf{证据分层}：A1–A11 是质量与配比实验数据，B4/B5 提供不同评测口径下的趋势对照，C1/C2/C4/C7/C8 提供榜单、模型与任务信息；B1 是近乎公式生成的经典锚点，B2、B6–B8 含半合成成分，B3 为插值，A12–A15 与 B10 为估算或外推，C3 为混合历史口径。拟合与检验分开：A1、A4/A5 和 B6 用于相应模型训练，A2/A3、A6–A11 与 B7 新增配置承担各自明确的检验角色；不能把异源 Loss 的绝对值视为同一交叉熵，也不能把合成数据上的拟合称为真实大模型验证。

\begin{figure}[H]
\centering\small
\renewcommand{\arraystretch}{1.12}
\begin{tabularx}{\textwidth}{@{}p{1.5cm}X p{4.2cm}@{}}
\toprule
阶段 & 建模对象与条件约束 & 新增量及向下游传递的结果 \\
\midrule
\textbf{问题一} & 文档质量、领域映射与 17 域配比；指标异质、部分域缺映射、配比份额和为 1 & $Q(x)\to Q_d\to Q(\bp)$；推荐配比 $\bp^*$ 与配比损失关系（P1） \\
\multicolumn{3}{c}{$\Downarrow$ 数据质量和配比进入规模模型} \\
\textbf{问题二} & 从 $L_0(N,D)$ 扩展质量项；A/B 评测尺度不同，B6 拟合、B7 新增配置留出 & 固定 $\bp^*$ 的 $L_{\mathrm{core}}(N,D,Q)$；配比通道暂为探索性（P2） \\
\multicolumn{3}{c}{$\Downarrow$ 标度律成为预算优化的目标} \\
\textbf{问题三} & 在预算 $C$、三项成本、质量边界和外生 $L_{\mathrm{ctx}}$ 下配置 $N,D,Q$ & $N^*,D^*,Q^*,L^*$ 与结构性转移状态（P3） \\
\multicolumn{3}{c}{$\Downarrow$ 最优损失与算力情景进入能力分析} \\
\textbf{问题四} & 在开放权重、模型类型、时间和评测口径约束下解释前沿 & 能力得分、规模与非规模模型贡献、条件预测（P4） \\
\bottomrule
\end{tabularx}
\caption{四问的递进关系：参数、约束与输出接口逐层进入下一问}
\label{fig:flow}
\end{figure}

\subsection{各问题的关键难点}

\textbf{问题一：从异质指标得到唯一语料质量。}目标是由 A1–A3 的文档信号建立 $Q(x)$，借 A16/A18 得到 17 域质量，再用 A4–A11 识别配比对 13 域损失的作用。难点在于 logits、规则量与 DSIR 的方向和尺度不一，长度与领域基线会干扰质量；17 域又有 11 域缺少直接评分。为避免冗余信号和极端指标支配结果，采用域内冻结预处理、CRITIC 权重与 Huber 中心；冲突只在评分后按来源诊断。领域质量以映射和文本桥接补齐，配比因份额闭合而采用中心对数比建模，并以 A6–A11 检查跨尺度边界。

\textbf{问题二：跨附件质量口径与可推广的函数形式。}目标是用 B1 的经典基线和 B6 的质量网格建立 $L(N,D,Q,\bp)$，并以 B7 新增配置、B4/B5 趋势及问题一接口检验。A、B 两组 Loss 的分词器与验证集可能不同，绝对值不能直接对齐。主口径取 $Q_{\mathrm B}=Q(\bp)$；质量乘子在嵌套族中比较，主式取幂次是为了写成有效规模，并同时报告留出误差更小的线性对照。配比效应保持独立，附件 A 只通过已经输出的 $Q$ 与 $\bp$ 进入。B6/B7 为半合成，限制了向真实训练的外推。

\textbf{问题三：非线性约束与边界转移。}目标是把前两问冻结的 $Q_0$、$\bp^*$ 和标度律参数作为输入，在给定 $C$ 与上下文长度下求 $N,D,Q$。训练、提质和注意力开销共同形成非线性约束，变量跨数量级，三种提质成本形状也不同。故先利用预算取等缩减维度，再分质量下界、内点和上界比较可行解，并以活跃约束变化判断结构性转移；不能仅凭单次局部优化或预设所有成本函数均凸。

\textbf{问题四：能力演进的解释边界。}目标是结合前问的规模状态与 C1/C2/C4/C6/C8 的能力、日期、算力和任务数据，解释前沿并作条件预测。难点是 Loss 与 Benchmark 得分尺度不同、任务会饱和，且规模与时间共变，C3 还有混合口径。故采用可检验的损失—得分桥接和连续时间项，在一致的开放权重与模型类型口径下分解增长，并作回测及窗口敏感性分析。时间项只表示控制规模后的剩余关联，不直接解释为已识别的因果技术进步。

\section{模型假设}

\begin{enumerate}[label=\arabic*.,leftmargin=2em]
\item \textbf{质量信号可作相对测量。}\label{ass:q1signal}A 组指标经方向与域内尺度校正后，可比较同域文档的相对质量；不假定它们有人工标注的绝对真值。服务问题一的 $Q(x)$；用 A2/A3 冻结评分与删去极端指标的扰动检查适用性。
\item \textbf{领域质量可按份额聚合。}\label{ass:q1mix}给定 A16 的映射与 A18 的桥接模型，混合语料质量取 17 域质量的份额加权值；配比对验证损失的额外影响单独建模。服务问题一至三；以六个可映射域的桥接误差、A6/A7 留出和 A8–A11 跨尺度直接检验检查边界。
\item \textbf{评价目标预先固定。}\label{ass:q1eval}13 个验证域等权定义主目标，pile\_cc 单域仅作敏感性口径，评价权重不随候选配比改变。服务问题一的配比搜索与问题三的目标选择；用两口径的候选差异检查敏感性。
\item \textbf{质量刻度需要锚定。}\label{ass:q2anchor}主口径取 $Q_{\mathrm B}=Q(\bp)$，推荐配比处 $Q_0=Q(\bp^*)=0.6624$。$Q=1$ 且 $\bp$ 等于 A4 训练配比均值时，广义标度律退回经典式。另一锚定 $Q_{\mathrm B}=Q(\bp)/Q_0$ 会使部分领域超出 $[0,1]$，只作敏感性。服务问题二、三；该锚点不能由 B1 的公式生成数据识别，B1 相对 Hoffmann 等（2022）常数的 RMSE 为 $1.5\times10^{-4}$。
\item \textbf{成本按同一训练量计。}\label{ass:q3cost}训练、提质和注意力开销均使用同一 $D$，成本函数参数依赛题给定。服务问题三；在三种函数和预算边界下分别检查最优解，不预设统一凸性。
\item \textbf{非规模变化在窗口内平滑。}\label{ass:q4time}问题四用连续时间项描述控制规模后的剩余得分变化，并以不同窗口与模型回测；遇到离散架构突破时，该条件预测不适用。
\end{enumerate}

\section{符号说明}

\begin{table}[H]
\centering\small
\caption{主要符号及含义}
\label{tab:symbol}
\begin{tabularx}{\textwidth}{@{}p{2.4cm}X p{1.55cm}p{2.6cm}@{}}
\toprule
符号 & 含义 & 单位 & 角色 \\
\midrule
$N$，$D$ & 模型参数量、训练 token 数 & 个，token & Q1/Q2 输入；Q3 决策变量 \\
$C$ & 总训练算力预算 & FLOPs & 给定参数（Q3） \\
$Q(x)$，$Q_i$ & 文档 $x$ 的主质量分、训练域 $i$ 的质量分 & $[0,1]$ & 估计量（Q1） \\
$Q(\bp)$，$Q_{\mathrm B}$ & 混合语料质量、B 组质量刻度；主口径二者相同 & 无量纲 & Q1 输出；Q2 输入 \\
$Q_0$ & 推荐配比下的提质起点，$Q(\bp^*)=0.6624$ & 无量纲 & 接口参数（Q2/Q3） \\
$\bp=(p_i)_{i=1}^{17}$ & 17 域训练份额，非负且总和为 1 & 无量纲 & 决策变量（Q1） \\
$z_j(x)$，$w_j$ & 文档第 $j$ 项标准化指标及冻结的 CRITIC 权重 & $[0,1]$，无量纲 & 数据变换、估计参数 \\
$B(x)$，$\tau_d$ & 六来源分歧度、A1 域 $d$ 的 P95 筛查阈值 & $[0,1]$ & 诊断量、固定阈值 \\
$L_v$，$L$ & 验证域 $v$ 的损失、等权综合损失 & nats & 观测量或预测目标 \\
$T_{iv}$ & 均匀配比附近，域 $i$ 份额对域 $v$ 的对数损失局部效应 & 无量纲 & 估计参数（Q1） \\
$E,A,B,\alpha,\beta$ & 经典标度律的地板、幅度与指数 & 依公式确定 & 估计参数（Q2） \\
$k_0,\kappa_N,\kappa_D$ & 质量地板及两项质量耦合强度 & nats，无量纲 & 估计参数（Q2） \\
$h_N(Q)$，$h_D(Q)$ & 主式质量乘子 $Q^{-\kappa_N}$、$Q^{-\kappa_D}$ & 无量纲 & 函数（Q2） \\
$s_{\mathrm{red}}$ & 配比乘在可约损失上的强度，1M 训练集定为 $1.500$ & 无量纲 & 估计参数（Q2） \\
$\varepsilon_N,\varepsilon_D,\varepsilon_Q$ & 损失对 $N,D,Q$ 的弹性 & 无量纲 & 导出量（Q2） \\
$g(Q)$，$L_{\mathrm{ctx}}$ & 单位 token 提质成本、上下文长度 & FLOPs/token，token & 给定函数与情景参数 \\
$C_{\mathrm{crit}}$ & 质量约束状态改变的临界预算 & FLOPs & 求解结果（Q3） \\
$S$，$t$ & Benchmark 综合得分、模型发布日期 & 分，年 & 响应量、时间变量（Q4） \\
\bottomrule
\end{tabularx}
\end{table}

表中只列全文反复使用的重要符号，局部符号在首次出现处说明。

\section{模型的建立与求解}

\subsection{数据来源与接口约定}

附件 A、B、C 分别承担质量与配比、标度律、能力演进的输入角色；同一附件内仍须区分拟合集、独立检验集和估算表。跨问只传递已定义的质量分、配比候选及模型参数，不以不同评测管线的 Loss 绝对值直接拼接。各项清洗、连接键、训练与留出划分在相应问题中说明；问题一的 A1–A3 指标处理见表~\ref{tab:q1preprocess}。
